\section{Combinatorics}

\subsection{Equally Likely Outcomes}
Now we consider some examples of probability spaces, where the outcomes are finite, and all outcomes are equally likely.

\begin{example}[Rolling a Fair Die]
    \label{ex:rolling-a-fair-die}%
    We consider the outcomes \(\Omega = \brce{1, 2, \ldots, 6}\), and \(\eventSpace = \powerset\para{\Omega}\) for all subsets of \(\Omega\) (since \(\Omega\) is countable).
        
    We naturally take for all \(\omega \in \Omega\),
    \[
        \Prob\para{\brce{\omega}} = \frac{1}{6},
    \]
    and for all \(A \subseteq \Omega\), we take
    \[
        \Prob\para{A} = \frac{\abs{A}}{6}.
    \]

    All outcomes are equally likely.
\end{example}

Lots of situations are like this in real life, and we refer to such setup as \emph{classical probability}, with \emph{equally likely outcomes}.

\begin{definition}[Equally Likely Outcomes]
    \label{def:equally-likely-outcomes}%
    For any finite set \(\Omega = \brce{\omega_1, \omega_2, \ldots, \omega_n}\), take \(\eventSpace = \powerset\para{\Omega}\). We define the probability measure to be
    \[
        \function{\Prob}{\eventSpace}{\ccintv{0}{1}}{A}{\frac{\abs{A}}{\abs{\Omega}}.}
    \]

    This is called the \emph{uniform probability distribution}.
\end{definition}

This models the experiment of picking a \emph{random} element of \(\Omega\). For all \(\omega \in \Omega\),
\[
    \Prob\para{\brce{\omega}} = \frac{1}{\abs{\Omega}}
\]

We look at some more examples on calculating the probability under such setup. Effectively, those are combinatorial problems counting the size of sets.

This example is more about counting the size of the \emph{sample space} which is the denominator.

\begin{example}[Picking Balls from a Bag]
    \label{ex:picking-balls-from-bag}%
    Suppose we have \(n\) balls labelled \(\brce{1, \ldots, n}\), and they are indistinguishable by touch. Pick \(k \leq n\) balls \emph{at random} (any outcome is equally likely) without replacement. We let the outcomes be
    \[
        \Omega = \set{A \subseteq \brce{1, 2, \ldots, n}}{\abs{A} = k}
    \]
    and hence
    \[
        \abs{\Omega} = \binom{n}{k}.
    \]

    Hence, for any \(\omega \in \Omega\), we have
    \[
        \Prob\para{\brce{\omega}} = \frac{1}{\binom{n}{k}}.
    \]
\end{example}

The following examples also require counting on the numerator, which is the cardinality of the event.

\begin{example}[Deck of Cards]
    \label{ex:deck-of-cards}%
    In a well-shuffled deck of 52 cards, take
    \[
        \Omega = \brce{\text{all permutations of 52 cards}}
    \]
    and hence
    \[
        \abs{\Omega} = 52! = 80,658,175,170,943,878,571,660,636,856,403,766,975,289,505,440,883,277,824,000,000,000,000.
    \]
    
    Take \(\eventSpace = \powerset\para{\Omega}\). Then
    \[
        \Prob\para{\text{top 2 cards are aces}} = \frac{4 \times 3 \times 50!}{52!}.
    \]
\end{example}

The next example uses results on set cardinalities.

\begin{example}[Largest Digit]
    \label{ex:largest-digit}%
    Consider a string of \(n\) random digits from \(0, 1, \ldots, 9\). Let
    \[
        \Omega = \brce{0, 1, \ldots, 9}^n,
    \]
    and we have \(\abs{\Omega} = 10^n\).

    Let the event \(A_k = \brce{\text{no digit exceeds }k}\), then
    \[
        \Prob\para{A_k} = \frac{\para{k + 1}^{n}}{10^n}.
    \]

    Now we consider the event \(B_k = \brce{\text{the largest digit is equal to }k}\), then note \(B_k = A_k \setminus A_{k - 1}\), so
    \[
        \Prob\para{B_k} = \Prob\para{A_k \setminus A_{k - 1}} = \Prob\para{A_k} - \Prob\para{A_{k - 1}} = \frac{\para{k + 1}^n - k^n}{10^n}.
    \]
\end{example}

\begin{example}[Birthday Problem]
    \label{ex:birthbay-problem}%
    There are \(n\) people. What is the probability that at least 2 of them share the same birthday?
        
    Assume each birthday is equally likely to be one of \(\brce{1, 2, \ldots, 365}\). Take
    \[
        \Omega = \brce{1, 2, \ldots, 365}^n,
    \]
    and take \(\eventSpace = \powerset\para{\Omega}\). Hence, for all \(\omega \in \Omega\), we have
    \[
        \Prob\para{\brce{\omega}} = \frac{1}{365^n}.
    \]

    Consider the event \(A = \brce{\text{at least 2 people have the same birthday}}\).

    Then \(A\compl = \brce{\text{all }n\text{ birthdays are different}}\). Hence,
    \[
        \Prob\para{A\compl} = \frac{365 \times 364 \times \cdots \times \para{365 - n + 1}}{365^n}
    \]
    and hence
    \[
        \Prob\para{A} = 1 - \Prob\para{A\compl}.
    \]

    Note that \(n = 22\) gives \(\Prob\para{A} \approx 0.476\), and \(n = 23\) gives \(\Prob\para{A} \approx 0.507\).
\end{example}

\subsection{Combinatorial Analysis}
We start off by investigating some easy problems in combinatorics.

\begin{example}[Partitioning Sets into Fixed Sizes]
    \label{ex:partition-sets-given-size}%
    We first consider \(\Omega\) being a finite set with \(\abs{\Omega} = n\). We want to partition \(\Omega\) into \(k\) disjoint subsets \(\Omega_1, \Omega_2, \ldots, \Omega_k\) with \(\Omega_i = n_i\) for \(\sum_{i = 1}^{k} n_i = n\).

    Let \(M\) be the number of ways to do this. We have
    \begin{align*}
        M &= \binom{n}{n_1} \times \binom{n - n_1}{n_2} \times \cdots \times \binom{n - \para{n_1 + \cdots + n_{k - 1}}}{n_k} \\
        &= \frac{n!}{n_1! n_2! \cdots n_k!}.
    \end{align*}
\end{example}

\begin{definition}[Multinomial Coefficient]
    \label{def:multinomial-coeff}%
    The \emph{multinomial coefficient} is defined as
    \[
        \binom{n}{n_1, n_2, \ldots, n_k} = \frac{n!}{n_1! n_2! \cdots n_k!}.
    \]
\end{definition}

\begin{examples}[Increasing and Strictly Increasing Functions]
    \label{ex:increasing-functions}%
    We say \(f\) is \emph{strictly increasing} if \(x < y\) implies \(f\para{x} < f\para{y}\), and \(f\) is \emph{(weakly) increasing}, if \(x \leq y\) implies \(f\para{x} \leq f\para{y}\).

    Let \(\functiondomain{f}{\brce{1, 2, \ldots, k}}{\brce{1, 2, \ldots, n}}\).
    
    \begin{enumerate}
        \item \textit{Strictly Increasing Functions.} Let \(k \leq n\). How many strictly increasing functions \(f\) are there? Such a function is uniquely determined by its range, so there are \(\binom{n}{k}\) of them.
        \item \textit{Increasing Functions.} How many increasing functions \(f\) are there?
        
        Let the function space
        \[
            \eventSpace\para{k, n} = \set{f}{\functiondomain{f}{\brce{1, 2, \ldots, k}}{\brce{1, 2, \ldots, n}}}.
        \]
        
        We define a bijection
        \[
            \set{f \in \eventSpace\para{k, n}}{f \text{ is increasing}} \to \set{g \in \eventSpace\para{k, n + k - 1}}{g \text{ is strictly increasing}},
        \]
        where \(g\para{i} = f\para{i} + i - 1\) that is strictly increasing. Hence, there are
        \[
            \binom{n + k - 1}{k}
        \]
        such increasing functions in \(\eventSpace\para{k, n}\).
    \end{enumerate}
\end{examples}

\subsection{Stirling's Formula}

The following asymptotic formula is very useful in the study of probability. Although the proof is non-examinable, we would still give a proof of it.

\begin{definition}[Asymptotic]
    \label{def:asymptotic}%
    For two sequences \(\para{a_n}, \para{b_n}\), we say \(\para{a_n}\) is \emph{asymptotic} to \(\para{b_n}\), denoted as
    \[
        a_n \sim b_n
    \]
    as \(n \to \infty\), if \(\frac{a_n}{b_n} \to 1\) as \(n \to \infty\).
\end{definition}

\begin{theorem}[Stirling's Formula]
    \label{thm:stirling-formula}%
    We have
    \[
        n! \sim n^n \sqrt{2 \pi n} \cdot e^{-n},
    \]
    as \(n \to \infty\).
\end{theorem}

This statement is relatively long to prove (and is non-examinable), but a statement that is also useful is given below, which is a corollary of the above statement, but whose proof is \textbf{examinable}.

\begin{corollary}
    \label{cor:weaker-statement-of-stirling}%
    We have \(\log \para{n!} \sim n \log n\) as \(n \to \infty\).
\end{corollary}

\begin{proof}
    Define
    \[
        l_n = \log\para{n!} = \log 2 + \cdots + \log n.
    \]

    For all \(x \in \RR\), denote \(\floor{x}\) as the integer part of \(x\). We have
    \[
        \log \floor{x} \leq \log x \leq \log \floor{x + 1}.
    \]

    We integrate this inequality from \(1\) to \(n\), and hence, we have
    \[
        \sum_{k = 1}^{n - 1} \log k \leq \int_{1}^{n} \log x \Diff x \leq \sum_{k = 1}^{n} \log k
    \]
    and hence we have
    \[
        l_{n - 1} \leq n \log n - n + 1 \leq l_n.
    \]

    Therefore, rearranging this gives
    \[
        n \log n - n + 1 \leq l_n \leq \para{n + 1} \log \para{n + 1} - \para{n + 1} + 1,
    \]
    and dividing this through by \(n \log n\) gives
    \[
        \frac{l_n}{n \log n} \to 1
    \]
    as \(n \to \infty\).
\end{proof}

We first show that Stirling's Formula is true, up to some constant (which would turn out to be \(\sqrt{2\pi}\)). A lemma that could be useful is a result that comes from integration by parts.

\begin{lemma}
    \label{lem:some-integration-formula}%
    For any \(\functiondomain{f}{\RR}{\RR}\) which is twice differentiable, for all \(a < b\), we have
    \[
        \int_a^b f\para{x} \Diff x = \frac{f\para{a} + f\para{b}}{2} \cdot \para{b - a} - \frac{1}{2} \int_a^b \para{x - a} \para{b - x} f''\para{x} \Diff x.
    \]
\end{lemma}

\begin{proof}
    Using integration by parts twice.
\end{proof}


\begin{lemma}[Stirling's Formula with Constant]
    \label{lem:stirling-formula-constant}%
    We have
    \[
        n! \sim A n^n \sqrt{n} \cdot e^{-n},
    \]
    as \(n \to \infty\), for some real number \(A \in \RR\).
\end{lemma}

\begin{proof}
    Take \(f\para{x} = \log x\), with \(a = k\), \(b = k + 1\), for \(k \in \NN\). Using \autoref{lem:some-integration-formula}, we have
    \[
        \int_{k}^{k + 1} \log x \Diff x = \frac{\log k + \log\para{k + 1}}{2} + \frac{1}{2} \int_{k}^{k + 1} \frac{\para{x - k} \para{k + 1 - x}}{x^2} \Diff x.
    \]

    Summing this from \(k = 1\) to \(n - 1\) gives
    \[
        \int_{1}^{n} \log x \Diff x = \frac{\log \para{n - 1}! + \log \para{n!}}{2} + \sum_{k = 1}^{n - 1} a_k,
    \]
    where
    \[
        a_k = \frac{1}{2} \int_{0}^{1} \frac{x \para{1 - x}}{\para{x + k}^2} \Diff x.
    \]

    Hence, we have
    \[
        n \log n - n + 1 = \log\para{n!} - \frac{\log n}{2} + \sum_{k = 1}^{n - 1} a_k
    \]

    Therefore,
    \[
        \log \para{n!} = n \log n - n + 1 + \frac{\log n}{2} - \sum_{k = 1}^{n - 1} a_k,
    \]
    and hence
    \[
        n! = n^n \cdot e^{-n} \cdot \sqrt{n} \cdot \exp\para{1 - \sum_{k = 1}^{n - 1} a_k}.
    \]

    Now, we have
    \[
        a_k = \frac{1}{2} \int_{0}^{1} \frac{x \para{1 - x}}{\para{x + k}^2} \Diff x \leq \frac{1}{2k^2} \int_{0}^{1} x\para{1 - x} \Diff x = \frac{1}{12k^2}.
    \]

    Therefore,
    \[
        \sum_{k = 1}^{\infty} a_k < \infty.
    \]

    Let \(A = \exp\para{1 - \sum_{k = 1}^{\infty} a_k}\). Therefore,
    \[
        n! = n^n \cdot e^{-n} \cdot \sqrt{n} \cdot A \cdot \exp\para{\sum_{k = n}^{\infty} a_k}.
    \]

    The \(\exp\para{\sum_{k = n}^{\infty} a_k}\) will converge to \(1\) as \(n \to \infty\), since the inner sum converges to \(0\).

    We have showed that
    \[
        \frac{n!}{n^n \cdot e^{-n} \cdot \sqrt{n}} \to A
    \]
    as \(n \to \infty\), as desired.
\end{proof}

To finish the proof of \autoref{thm:stirling-formula} (Stirling's Formula), we want to find \(A\). An asymptotic analysis on the binomial coefficient \(2^{-2n} \binom{2n}{n}\) would be useful, and one method to do this is to use the expansion of the binomial coefficient.

\begin{lemma}[Asymptotic for \(2^{-2n} \binom{2n}{n}\), using Factorials]
    \label{lem:asymptotic-binomial-factorial}%
    We have
    \[
        2^{-2n} \binom{2n}{n} \sim \frac{\sqrt{2}}{A \sqrt{n}}
    \]
    for the \(A\) above, in \autoref{lem:stirling-formula-constant}.
\end{lemma}

\begin{proof}
    Notice that
    \begin{align*}
        2^{-2n} \binom{2n}{n} &= 2^{-2n} \cdot \frac{\para{2n}!}{n! \cdot n!} \\
        &\sim \frac{2^{-2n} \cdot \para{2n}^{2n} \cdot e^{-2n} \cdot \sqrt{2n} \cdot A}{n^n \cdot e^{-n} \cdot \sqrt{n} \cdot A \cdot n^n \cdot e^{-n} \cdot \sqrt{n} \cdot A} \\
        &= \frac{\sqrt{2}}{A \sqrt{n}}.
    \end{align*}
\end{proof}

Another way to study the asymptotic behaviour of this function, is to relate this with the Wallis Integral.

\begin{lemma}[The Wallis Integrals]
    \label{lem:wallis-integral}%
    Let
    \[
        I_n = \int_{0}^{\frac{\pi}{2}} \cos^n \theta \Diff \theta
    \]
    be the Wallis Integrals. We have
    \[
        I_{2n} = \frac{\pi}{2} \cdot 2^{-2n} \cdot \binom{2n}{n},
    \]
    and
    \[
        I_{2n + 1} = \frac{1}{2n + 1} \cdot \para{2^{-2n} \cdot \binom{2n}{n}}\inverse.
    \]
\end{lemma}

\begin{proof}
    We notice \(I_0 = \frac{\pi}{2}\) and \(I_2 = 1\). Integration by parts gives us
    \[
        I_n = \frac{n - 1}{n} \cdot I_{n - 2},
    \]
    so
    \begin{align*}
        I_{2n} &= \frac{2n - 1}{2n} \cdot I_{2n - 2} \\
        &= \cdots \\
        &= \frac{\para{2n - 1} \cdot \para{2n - 3} \cdots 3 \cdot 1}{\para{2n} \cdot \para{2n - 2} \cdots 2} I_0 \\
        &= \frac{\para{2n} \cdot \para{2n - 1} \cdot \para{2n - 2} \cdot \para{2n - 3} \cdots 3 \cdot 2 \cdot 1}{2n \cdot 2n \cdot \para{2n - 2} \cdot \para{2n - 2} \cdots 2 \cdot 2} \cdot \frac{\pi}{2} \\
        &= \frac{\para{2n}!}{2^{2n} \cdot n! \cdot n!} \cdot \frac{\pi}{2} \\
        &= \frac{\pi}{2} \cdot 2^{-2n} \cdot \binom{2n}{n},
    \end{align*}
    so we have
    \[
        I_{2n} = \frac{\pi}{2} \cdot 2^{-2n} \cdot \binom{2n}{n}.
    \]

    Similarly, we have
    \[
        I_{2n + 1} = \frac{2n \cdot \para{2n - 1} \cdots 4 \cdot 2}{\para{2n + 1} \cdot \para{2n - 1} \cdots 3 \cdot 1} I_1 = \frac{1}{2n + 1} \cdot \para{2^{-2n} \cdot \binom{2n}{n}}\inverse.
    \]

\end{proof}

\begin{corollary}[Asymptotic for \(2^{-2n} \binom{2n}{n}\), using Wallis]
    \label{cor:asymptotic-binomial-wallis}%
    We have
    \[
        2^{-2n} \binom{2n}{n} \sim \frac{1}{\sqrt{\pi n}}.
    \]
\end{corollary}

\begin{proof}
    We use \autoref{lem:wallis-integral}. 

    Notice that we have \(\frac{I_n}{I_{n - 2}} \to 1\) as \(n \to \infty\).

    Note that \(I_n\) is decreasing in \(n\), so
    \[
        \frac{I_{2n}}{I_{2n + 1}} \leq \frac{I_{2n - 1}}{I_{2n + 1}} \to 1
    \]
    and
    \[
        \frac{I_{2n}}{I_{2n + 1}} \geq \frac{I_{2n}}{I_{2n - 1}} \to 1,
    \]
    so we have \(\frac{I_{2n}}{I_{2n + 1}} \to 1\) as \(n \to \infty\).

    Hence,
    \[
        \frac{\frac{\pi}{2} \cdot 2^{-2n} \cdot \binom{2n}{n}}{\frac{1}{2n + 1} \cdot \para{2^{-2n} \cdot \binom{2n}{n}}\inverse} \to 1
    \]
    which gives
    \[
        \para{2^{-2n} \cdot \binom{2n}{n}}^2 \sim \frac{1}{\pi n}
    \]
    as \(n \to \infty\), as desired.
\end{proof}

Now we are ready to find \(A\), and give a proof of the original statement in \autoref{thm:stirling-formula} (Stirling's formula).

\begin{proof}[of \autoref{thm:stirling-formula} (Stirling's Formula)]
    From \autoref{lem:stirling-formula-constant}, it remains to show that \(A = \sqrt{2 \pi}\).
    
    Using the result in \autoref{lem:asymptotic-binomial-factorial}, we have
    \[
        2^{-2n} \binom{2n}{n} \sim \frac{\sqrt{2}}{A \sqrt{n}}.
    \]

    From \autoref{cor:asymptotic-binomial-wallis}, we also have
    \[
        2^{-2n} \binom{2n}{n} \sim \frac{1}{\sqrt{\pi n}}.
    \]
    
    This then implies \(A = \sqrt{2 \pi}\), which shows the result as desired.
\end{proof}

